\documentclass[12pt]{article}
\usepackage{latexblog}

% ============================================================
% Blog Metadata
% ============================================================
\blogtitle{C++中的NRVO}
\blogdate{2019-09-25}
\blogtags{编程语言, 闲话编程}
\bloglang{zh}
\blogsummary{}

\begin{document}
\makeblogtitle

对于C++这种需要精细管理对象的语言来说有时候真是比较复杂，一个看似简单的问题一直在困惑着我：到底可不可以在方法中返回局部变量呢？

\section{可以返回临时变量}\label{ux53efux4ee5ux8fd4ux56deux4e34ux65f6ux53d8ux91cf}

答案是肯定的，如果我们在一个方法中返回了临时变量，这个临时变量实际上是在栈里面的，当执行完方法后栈就销毁了，那么为什么我们还可以这样做呢？来看一个例子：

\begin{Shaded}
\begin{Highlighting}[]
\PreprocessorTok{\#include }\ImportTok{\textless{}iostream\textgreater{}}
\KeywordTok{using} \KeywordTok{namespace}\NormalTok{ std}\OperatorTok{;}

\KeywordTok{class}\NormalTok{ Value }\OperatorTok{\{}
    \KeywordTok{public}\OperatorTok{:}
\NormalTok{        Value}\OperatorTok{(}\DataTypeTok{int}\NormalTok{ \_m}\OperatorTok{):}\NormalTok{m}\OperatorTok{(}\NormalTok{\_m}\OperatorTok{)} \OperatorTok{\{} \BuiltInTok{std::}\NormalTok{cout }\OperatorTok{\textless{}\textless{}} \StringTok{"test constructor"} \OperatorTok{\textless{}\textless{}}\NormalTok{ m }\OperatorTok{\textless{}\textless{}} \BuiltInTok{std::}\NormalTok{endl}\OperatorTok{;} \OperatorTok{\}}
\NormalTok{        Value}\OperatorTok{(}\AttributeTok{const}\NormalTok{ Value}\OperatorTok{\&}\NormalTok{ t}\OperatorTok{)} \OperatorTok{\{} 
            \BuiltInTok{std::}\NormalTok{cout }\OperatorTok{\textless{}\textless{}} \StringTok{"test copy constructor"} \OperatorTok{\textless{}\textless{}}\NormalTok{ m }\OperatorTok{\textless{}\textless{}} \BuiltInTok{std::}\NormalTok{endl}\OperatorTok{;} 
            \KeywordTok{this}\OperatorTok{{-}\textgreater{}}\NormalTok{m }\OperatorTok{=}\NormalTok{ t}\OperatorTok{.}\NormalTok{m}\OperatorTok{;}
        \OperatorTok{\}}
        \OperatorTok{\textasciitilde{}}\NormalTok{Value}\OperatorTok{()} \OperatorTok{\{} \BuiltInTok{std::}\NormalTok{cout }\OperatorTok{\textless{}\textless{}} \StringTok{"test destructor"} \OperatorTok{\textless{}\textless{}}\NormalTok{ m }\OperatorTok{\textless{}\textless{}} \BuiltInTok{std::}\NormalTok{endl}\OperatorTok{;} \OperatorTok{\}}
        \DataTypeTok{void}\NormalTok{ print}\OperatorTok{()} \OperatorTok{\{} \BuiltInTok{std::}\NormalTok{cout }\OperatorTok{\textless{}\textless{}} \StringTok{"m:"} \OperatorTok{\textless{}\textless{}}\NormalTok{ m }\OperatorTok{\textless{}\textless{}} \BuiltInTok{std::}\NormalTok{endl}\OperatorTok{;} \OperatorTok{\}}
    \KeywordTok{private}\OperatorTok{:}
        \DataTypeTok{int}\NormalTok{ m}\OperatorTok{;}
\OperatorTok{\};}

\KeywordTok{class}\NormalTok{ Producer }\OperatorTok{\{}
\KeywordTok{public}\OperatorTok{:}
\NormalTok{    Value produce}\OperatorTok{(}\DataTypeTok{int}\NormalTok{ i}\OperatorTok{)} \OperatorTok{\{}
\NormalTok{        Value t}\OperatorTok{(}\NormalTok{i}\OperatorTok{);}
        \ControlFlowTok{return}\NormalTok{ t}\OperatorTok{;}
    \OperatorTok{\}}
\OperatorTok{\};}

\DataTypeTok{int}\NormalTok{ main}\OperatorTok{(}\DataTypeTok{int}\NormalTok{ argc}\OperatorTok{,} \DataTypeTok{char}\OperatorTok{*}\NormalTok{ argv}\OperatorTok{[])} \OperatorTok{\{}
    \BuiltInTok{std::}\NormalTok{cout }\OperatorTok{\textless{}\textless{}} \StringTok{"Hello world!"} \OperatorTok{\textless{}\textless{}} \BuiltInTok{std::}\NormalTok{endl}\OperatorTok{;}
\NormalTok{    Producer p}\OperatorTok{;}
\NormalTok{    Value t }\OperatorTok{=}\NormalTok{ p}\OperatorTok{.}\NormalTok{produce}\OperatorTok{(}\DecValTok{100}\OperatorTok{);}
\NormalTok{    t}\OperatorTok{.}\NormalTok{print}\OperatorTok{();}
    \ControlFlowTok{return} \DecValTok{1}\OperatorTok{;}
\OperatorTok{\}}
\end{Highlighting}
\end{Shaded}

执行的结果是：

\begin{verbatim}
Hello world!
test constructor100
m:100
test destructor100
\end{verbatim}

结果证明这样做其实是可以取到我们定义的值的，这么做可行的原因是，实际上，编译器会帮我们把临时变量拷贝一份出来，所以即便栈销毁了，我们也能够拿到新的值。

\section{不要返回临时变量的引用}\label{ux4e0dux8981ux8fd4ux56deux4e34ux65f6ux53d8ux91cfux7684ux5f15ux7528}

那么，如果我们返回临时变量的引用呢？

\begin{Shaded}
\begin{Highlighting}[]
\KeywordTok{class}\NormalTok{ Producer }\OperatorTok{\{}
\KeywordTok{public}\OperatorTok{:}
\NormalTok{    Value}\OperatorTok{*}\NormalTok{ produce}\OperatorTok{(}\DataTypeTok{int}\NormalTok{ i}\OperatorTok{)\{}
\NormalTok{        Value t}\OperatorTok{(}\NormalTok{i}\OperatorTok{);}
        \ControlFlowTok{return} \OperatorTok{\&}\NormalTok{t}\OperatorTok{;}
    \OperatorTok{\}}
\OperatorTok{\};}
\end{Highlighting}
\end{Shaded}

这样做得到的结果是不对的：

\begin{verbatim}
test.cpp:21:17: warning: address of stack memory associated with local variable 't' returned [-Wreturn-stack-address]
        return &t;
                ^
1 warning generated.
Hello world!
test constructor100
test destructor100
m:-327065280
\end{verbatim}

编译器会有一个警告，尽管我们仍旧可以运行我们的代码，但是实际上我们得到的值是不对的。

\section{NRVO机制}\label{nrvoux673aux5236}

那么，既然我们返回临时对象的值，实际上会得到一个拷贝的对象，那么如果我们有拷贝构造函数，是不是应该被调用呢？

然而在前面的例子中，拷贝构造函数并没有被调用到，这又是为什么呢？答案就是因为NRVO(Return Value Optimization)。这是c++11中的特性。我们首先可以尝试禁用掉这个特性，看看会发生什么。

\begin{verbatim}
hfli@CNhfli ~ $ g++ -fno-elide-constructors test.cpp
hfli@CNhfli ~ $ ./a.out
Hello world!
test constructor100
test copy constructor0
test destructor100
test copy constructor0
test destructor100
m:100
test destructor100
\end{verbatim}

可以看出，拷贝构造函数调用了两次，第一次是在produce函数中返回的时候，第二次是我们在给变量赋值的时候。


\end{document}
